\documentclass[../../main.tex]{subfiles}
\begin{document}
% SOURCE p=240; printed=225; object=ch07-ex20-czero-nonreflexive; number=20; region=upper
\begin{exercise}\label{exercise:ch07-czero-nonreflexive}
下面给出 $C_0(X)$ 不自反的一些例子。
\begin{enumerate}
\item\label{item:ch07-atom-sum-functional} 若 $\mu\in M(X)$，令 $\Phi(\mu)=\sum_{x\in X}\mu(\{x\})$。该和有良好定义，且 $\Phi\in M(X)^*$。若存在非零 $\mu\in M(X)$ 使所有 $x\in X$ 都满足 $\mu(\{x\})=0$，则 $\Phi$ 不在 $C_0(X)$ 到 $M(X)^*\cong C_0(X)^{**}$ 的典范映射的像中。
\item\label{item:ch07-discrete-czero-dual} 在另一个极端，令 $X=\mathbb N$ 配备离散拓扑，则 $C_0(X)^*\cong\ell^1$，且 $(\ell^1)^*\cong\ell^\infty$。（注意：$C_0(\mathbb N)$ 通常记为 $c_0$。）
\end{enumerate}
\end{exercise}
% SUPPLEMENT: id=v14-ch07-ex20-solution; parent=ch07-ex20-czero-nonreflexive; kind=solution+motivation+hypotheses; origin=independent-derivation
\begin{solution}
\textbf{逐项补写（重构新增）。}第一项：对每个有限子集 $F\subseteq X$，单点互不相交，故 $\sum_{x\in F}|\mu(\{x\})|\leqslant|\mu|(X)$。因此非负无序和有限；特别地，对每个 $n$，满足 $|\mu(\{x\})|\geqslant1/n$ 的点只能有限个，非零项至多可数。复数和绝对收敛，且 $|\Phi(\mu)|\leqslant\|\mu\|$。对两测度，可在其非零原子集合的可数并上使用绝对收敛级数的线性性，得到 $\Phi$ 线性，所以它有界。

假设存在 $f\in C_0(X)$，使 $\Phi(\lambda)=\int f\,\mathrm d\lambda$ 对所有 $\lambda\in M(X)$ 成立。每个 $x$ 的点质量 $\delta_x$ 为 Radon 测度：若 Borel 集包含 $x$，紧集 $\{x\}$ 给出内逼近；若不包含 $x$，开集 $X\setminus\{x\}$ 给出零质量外罩。因此代入 $\delta_x$ 得 $f(x)=1$。若 $X$ 非紧，常数 $1$ 不属于 $C_0(X)$，已经矛盾。

无论 $X$ 是否紧，再用题设的非零无原子测度 $\mu$：存在 Borel 集 $E$ 使 $\mu(E)\ne0$。令 $\nu(B)=\mu(B\cap E)$，则 $|\nu|=\chi_E|\mu|$；由可积权函数保持 Radon 性，$\nu\in M(X)$。每个单点仍有 $\nu(\{x\})=0$，所以 $\Phi(\nu)=0$；但 $f\equiv1$ 又给 $\int f\,\mathrm d\nu=\mu(E)\ne0$，矛盾。这里特意使用 $E$，因为非零复测度本身可能满足 $\mu(X)=0$。

第二项：离散 $\mathbb N$ 的紧集恰为有限集，所以 $C_0(\mathbb N)=c_0$。每个有限复测度由 $a_n=\mu(\{n\})$ 决定，$\mu(E)=\sum_{n\in E}a_n$，且 $\|\mu\|=\sum_n|a_n|$。反向每个 $a\in\ell^1$ 都定义这种测度；它在所有集合上由有限集内逼近，且所有集合本身开，故 Radon。因此 Riesz 定理给 $c_0^*\cong M(\mathbb N)\cong\ell^1$，等距同构。

对 $b\in\ell^\infty$，$L_b(a)=\sum_na_nb_n$ 为 $\ell^1$ 上有界线性泛函，$\|L_b\|=\|b\|_\infty$；上界由绝对和，反向由单坐标单位向量取上确界得到。对任意 $L\in(\ell^1)^*$，置 $b_n=L(e_n)$，则 $|b_n|\leqslant\|L\|$。对有限序列线性性给 $L=L_b$，再由有限截断在 $\ell^1$ 中稠密推广到所有序列。因此 $(\ell^1)^*=\ell^\infty$。典范嵌入对应 $c_0\subsetneq\ell^\infty$，例如常序列 $1$ 不在 $c_0$，故 $c_0$ 不自反。
\end{solution}
\begin{remark}
\textbf{动机与条件（重构新增）。}对测度取“全部原子的总量”是一个有界操作，却不能总由一个连续测试函数完成。非零无原子测度使这种差别可检测；有限总变差则保证无序和绝对收敛。第二项说明即使所有测度都由原子组成，双对偶仍可能增加不在无穷远消失的函数；因此没有无原子测度并不推出自反。
\end{remark}
% END SUPPLEMENT: id=v14-ch07-ex20-solution
% SOURCE p=240; printed=225; object=ch07-ex21-finite-moment-constraints; number=21; region=upper-middle
\begin{exercise}\label{exercise:ch07-finite-moment-constraints}
设 $\{f_\alpha\}_{\alpha\in A}\subseteq C_0(X)$，$\{c_\alpha\}_{\alpha\in A}$ 为复数族。若对每个有限集 $B\subseteq A$，存在 $\mu_B\in M(X)$，使 $\|\mu_B\|\leqslant1$ 且对 $\alpha\in B$ 有 $\int f_\alpha\,\mathrm d\mu_B=c_\alpha$，则存在 $\mu\in M(X)$，使 $\|\mu\|\leqslant1$ 且对全部 $\alpha\in A$ 有 $\int f_\alpha\,\mathrm d\mu=c_\alpha$。
\end{exercise}
% SUPPLEMENT: id=v14-ch07-ex21-solution; parent=ch07-ex21-finite-moment-constraints; kind=solution+motivation+hypotheses; origin=independent-derivation
\begin{solution}
\textbf{完整补写（重构新增）。}在 $V=\operatorname{span}\{f_\alpha:\alpha\in A\}$ 上尝试定义 $L(\sum_{j=1}^ma_jf_{\alpha_j})=\sum_{j=1}^ma_jc_{\alpha_j}$。若两个有限表示给出同一函数，把两组下标合在有限集 $B$ 中；用 $\mu_B$ 积分其差，得到对应的系数加权 $c_\alpha$ 之差为零，因此定义与表示无关。线性性随之成立。对 $v=\sum_ja_jf_{\alpha_j}$ 选包含这些下标的 $B$，有
\begin{align*}
|L(v)|=\left|\int v\,\mathrm d\mu_B\right|\leqslant\|v\|_u\|\mu_B\|\leqslant\|v\|_u.
\end{align*}
Hahn--Banach 定理把 $L$ 延拓为 $C_0(X)$ 上范数至多 $1$ 的泛函，Riesz 定理将其表示为所需 $\mu$。若 $A$ 为空，取 $\mu=0$ 即可。
\end{solution}
\begin{remark}
\textbf{动机与条件（重构新增）。}任意有限组相容的积分读数在统一范数约束下可拼成全体相容的读数。统一的 $1$ 控制使有限维数据在其线性张成空间上有同一连续界，Hahn--Banach 才能保留这个界；可把 $1$ 换成任何固定有限常数。完全不要求指标集可数，所以证明没有把任意族误当作一个序列。
\end{remark}
% END SUPPLEMENT: id=v14-ch07-ex21-solution
% SOURCE p=240; printed=225; object=ch07-ex22-czero-weak-sequences; number=22; region=middle
\begin{exercise}\label{exercise:ch07-czero-weak-sequences}
$C_0(X)$ 中的序列 $\{f_n\}$ 弱收敛到 $f\in C_0(X)$，当且仅当 $\sup_n\|f_n\|_u<\infty$ 且 $f_n\to f$ 逐点成立。
\end{exercise}
% SUPPLEMENT: id=v14-ch07-ex22-solution; parent=ch07-ex22-czero-weak-sequences; kind=solution+motivation+hypotheses; origin=independent-derivation
\begin{solution}
\textbf{完整补写（重构新增）。}若弱收敛，则用 $\delta_x\in C_0(X)^*$ 测试得到 $f_n(x)\to f(x)$。把 $f_n$ 看作双对偶中的泛函 $Jf_n$，对每个 $I\in C_0(X)^*$，数列 $(Jf_n)(I)=I(f_n)$ 收敛，故有界。对 Banach 空间 $C_0(X)^*$ 使用一致有界原理，得到 $\sup_n\|Jf_n\|<\infty$；典范嵌入等距，故 $\sup_n\|f_n\|_u<\infty$。

反向设一致界为 $M$ 且逐点趋于已经属于 $C_0(X)$ 的 $f$，则 $|f|\leqslant M$。任给 $I\in C_0(X)^*$，用 Riesz 定理写成 $I_\mu$。$|\mu|$ 为有限正测度，且 $|f_n-f|\leqslant2M$，控制收敛给
\begin{align*}
|I(f_n)-I(f)|\leqslant\int|f_n-f|\,\mathrm d|\mu|\longrightarrow0.
\end{align*}
因此 $f_n\rightharpoonup f$。
\end{solution}
\begin{remark}
\textbf{动机与条件（重构新增）。}在有限测度下，一致有界的逐点收敛足以控制全部积分测试；Riesz 定理恰把这些测试穷尽。极限预先属于 $C_0$ 是条件，不能从逐点收敛自动推出连续或在无穷远消失。这里使用的是序列控制收敛，不能未经额外论证推广为任意网的同一判据。
\end{remark}
% END SUPPLEMENT: id=v14-ch07-ex22-solution
% SOURCE p=240; printed=225; object=ch07-ex23-distribution-assumptions; number=23; region=middle
\begin{exercise}\label{exercise:ch07-distribution-assumptions}
\subcref{proposition:ch07-vague-distribution}{item:ch07-vague-implies-distribution}中的假设 $\mu_n\geqslant0$ 及 $\lim_{x\to-\infty}[\sup_nF_n(x)]=0$ 不能省略。（分别取 $\mu_n$ 为 $n^{-1}$ 与 $-n^{-1}$ 处点质量之差，或 $-n$ 处的点质量。）
\end{exercise}
% SUPPLEMENT: id=v14-ch07-ex23-solution; parent=ch07-ex23-distribution-assumptions; kind=solution+motivation+hypotheses; origin=independent-derivation
\begin{solution}
\textbf{两个反例（重构新增）。}先取 $\mu_n=\delta_{1/n}-\delta_{-1/n}$，$\mu=0$。对 $f\in C_0(\mathbb R)$，$\int f\,\mathrm d\mu_n=f(1/n)-f(-1/n)\to0$，故模糊收敛。所有测度都支撑于 $[-1,1]$，因此 $x<-1$ 时全部 $F_n(x)=0$，统一左尾条件满足。然而 $F_n(0)=-1$，而 $F\equiv0$ 在 $0$ 连续，所以分布函数不收敛。这里缺少的恰为正性。

再取 $\mu_n=\delta_{-n}\geqslant0$、$\mu=0$。由于 $f\in C_0$ 在负无穷处趋于零，$\int f\,\mathrm d\mu_n=f(-n)\to0$。但对每个固定 $x$，当 $n$ 足够大时 $-n\leqslant x$，所以 $F_n(x)=1\to1\ne F(x)=0$。对每个 $x$ 又有 $\sup_nF_n(x)=1$，故恰是统一左尾条件失效。
\end{solution}
\begin{remark}
\textbf{条件的含义（重构新增）。}“不能省略”指删去任一条件后一般结论不再成立，并非声称每个恰好收敛的特例都必须满足这些充分条件。第一个例子展示正负质量在连续测试下抵消，第二个展示连续消失测试看不见逃向负无穷的质量。
\end{remark}
% END SUPPLEMENT: id=v14-ch07-ex23-solution
% SOURCE p=240; printed=225; object=ch07-ex24-vague-counterexamples; number=24; region=middle-lower
\begin{exercise}\label{exercise:ch07-vague-counterexamples}
找出 $M(\mathbb R)$ 中满足下列条件的序列 $\{\mu_n\}$。
\begin{enumerate}
\item\label{item:ch07-vague-not-norm} $\mu_n\to0$ 模糊收敛，但 $\|\mu_n\|\not\to0$。
\item\label{item:ch07-vague-not-measurable-test} $\mu_n\to0$ 模糊收敛，但对某个具有紧支集的有界可测 $f$，$\int f\,\mathrm d\mu_n\not\to\int f\,\mathrm d\mu$。
\item\label{item:ch07-vague-not-distribution} $\mu_n\geqslant0$ 且 $\mu_n\to0$ 模糊收敛，但存在 $x\in\mathbb R$ 使 $F_n(x)\not\to F(x)$（记号同\cref{proposition:ch07-vague-distribution}）。
\end{enumerate}
\end{exercise}
% SUPPLEMENT: id=v14-ch07-ex24-solution; parent=ch07-ex24-vague-counterexamples; kind=solution+motivation+hypotheses; origin=independent-derivation
\begin{solution}
\textbf{逐项补写（重构新增）。}各项中的极限测度均为 $\mu=0$。
\begin{enumerate}
\item 取 $\mu_n=\delta_n$。每个 $f\in C_0(\mathbb R)$ 满足 $f(n)\to0$，所以模糊收敛，但 $\|\mu_n\|=1$。
\item 取 $\mu_n=\delta_{1/n}-\delta_0$。对连续 $g$，$g(1/n)-g(0)\to0$，所以模糊收敛。取 $f=\chi_{\{0\}}$，它 Borel 可测、有界，支集为紧集 $\{0\}$，但 $\int f\,\mathrm d\mu_n=-1$，而 $\int f\,\mathrm d\mu=0$。
\item 取 $\mu_n=\delta_{-n}$，则为非零正测度，且上一题已证明其模糊趋于零。在 $x=0$ 处，$F_n(0)=1$ 而 $F(0)=0$。
\end{enumerate}
\end{solution}
\begin{remark}
\textbf{动机与条件（重构新增）。}三例依次区分模糊收敛与范数收敛、连续测试与可测测试、消失测试与半直线测试。第二项说明仅要求测试函数有界且紧支撑，仍无法替代连续性。第三项不与收敛判据矛盾，因为缺少统一左尾控制。
\end{remark}
% END SUPPLEMENT: id=v14-ch07-ex24-solution
\subfile{section-03-exercises-approximation.tex}
\end{document}
